\tikzset{add/.style n args={4}{
    minimum width=14,
    path picture={
    \draw[black] 
        (path picture bounding box.south east) -- (path picture bounding box.north west)
        (path picture bounding box.south west) -- (path picture bounding box.north east);
    \node[anchor=north, inner sep=0pt] at ($(path picture bounding box.north)+(0,-0.13)$) {\tiny #1};
    \node[anchor=east,  inner sep=0pt] at ($(path picture bounding box.east)+(-0.13,0)$)  {\tiny #2};
    \node[anchor=south, inner sep=0pt] at ($(path picture bounding box.south)+(0,0.13)$) {\tiny #3};
    \node[anchor=west,  inner sep=0pt] at ($(path picture bounding box.west)+(0.13,0)$)  {\tiny #4};
        }
    }
}

\tikzstyle{block} = [draw, rectangle, align=center, font=\tiny]%
\tikzstyle{sum} = [draw, circle, add={}{+}{+}{}, font=\tiny]%
\tikzstyle{comp} = [draw, circle, add={}{}{+}{$-$}, font=\tiny]%
\tikzstyle{input} = [text centered, inner sep=0pt, outer sep=0pt, font=\footnotesize]%
\tikzstyle{output} = [text centered, inner sep=0pt, outer sep=0pt, font=\footnotesize]%

\begin{tikzpicture}[auto,node distance=1cm,>=latex']
    % ---------------------------------
	% --- taille de font par défaut ---
	% ---------------------------------
    \scriptsize

	% ----------------------------------
	% --- définition des espacements ---
	% ----------------------------------
	\def\x{1em}
	\def\y{2em}
    \def\dx{-470}
    \def\dy{1}
    \def\iv{8em}

    % ----------------------------
	% --- définition des nodes ---
	% ----------------------------
	% node de consigne
	\node [input] (r2) {$r_2(k)$};

    % node des régulateurs + préfiltre
    \node [block, below=2*\y of r2.south, fill=cyan!30, anchor=north] (CpCiPf) {$C_P(q)+C_I(q)P\!F(q)$}; 

    % node sommateur
    \node [sum, below=\y of CpCiPf.south, fill=cyan!30, anchor=north] (sum2) {};

    % node comparateur
    \node [comp, below=\y of sum2.south, fill=cyan!30, anchor=north] (comp2) {};

    % node des régulateurs
    \node [block, right=\x of comp2.east, fill=cyan!30, anchor=west] (CpCi) {$C_P(q)+C_I(q)$}; 

    % node de commande
    \node [input] (r1) at ([xshift=-\x]r2-|CpCiPf.west) {$r_1(k)$};

    % node de sortie
    \node [input] (y) at ([xshift=\x]r2-|CpCi.east) {$y(k)$};

    % node filtres
    \node [block, below=\y of comp2.south, fill=yellow!30, anchor=north] (Fsvfl) {$A^{-1}(q,\hat{\theta}^{\text{iter}-1}_{sriv})$};
    \node [block, fill=yellow!30] (Fsvfr) at (Fsvfl-|y) {$A^{-1}(q,\hat{\theta}^{\text{iter}-1}_{sriv})$};

    % node vecteur régresseur
    \node [block, fill=cyan!30, anchor=north] (phi) at ([yshift=-\y]$(Fsvfl.south)!0.5!(Fsvfr.south)$){%
        $\begin{aligned}%
              \varphi_f^T(k) = [ - & y_f(k-n_a),\,\cdots, -y_f(k-1),                       \\%
                             & u_f(k-n_b-n_k-1), \,\cdots, \,u_f(k-n_k) ]%
          \end{aligned}$%
    };

    % node matrice de régression
    \node [block, below=\y of phi.south, fill=cyan!30, anchor=north] (Phi) {%
        $\Phi_f = %
        \begin{bmatrix} %
            \varphi_f^T(1) \\ %
            \vdots \\ %
            \varphi_f^T(N) %
        \end{bmatrix}$%
    };

    % node matrice de sortie
    \node [block, below right=\y and \x of phi.south east, fill=cyan!30, anchor=north] (Y) {%
        $Y_f = %
        \begin{bmatrix} %
            y_f(1) \\ %
            \vdots \\ %
            y_f(N) %
        \end{bmatrix}$%
    };

    % node solution
    \node [block, below=\y of Phi.south, fill=cyan!30, anchor=north] (sol) {\textbf{Solution CLSRIVFR}\\$\hat{\theta}_{sriv}=\left(Z_f^T \Phi_f\right)^{-1}Z_f^T Y_f$};

    % node de sortie
    \node [output, below=2*\y of sol.south, anchor=north] (theta) {$\hat{\theta}_{sriv}$};

    % node fonction de sensibilité auxiliaire
    \node [block, left=\iv of CpCiPf.west, fill=teal!30, anchor=east] (Saux) {$\displaystyle \frac{1}{1+\left[C_P(q)+C_I(q)\right]\hat{G}^{\text{iter}}\left(q,\hat{\theta}_{sriv}^{\text{iter}-1}\right)}$};

    % node système auxiliaire
    \node [block, below=\y of Saux.south, fill=teal!30, anchor=north] (Gaux) {$\hat{G}^{\text{iter}}\left(q,\hat{\theta}_{sriv}^{\text{iter}-1}\right)$};

    % node filtres auxiliaires
    \node [block, fill=yellow!30] (Fsvflaux) at (Gaux|-Fsvfl) {$A^{-1}(q,\hat{\theta}^{\text{iter}-1}_{sriv})$};
    \node [block, fill=yellow!30, anchor=west] (Fsvfraux) at ([xshift=\x]Fsvflaux.east) {$A^{-1}(q,\hat{\theta}^{\text{iter}-1}_{sriv})$};

    % node vecteur instrumental
    \node [block, fill=teal!30, anchor=north] (z) at ([yshift=-\y]$(Fsvflaux.south)!0.5!(Fsvfraux.south)$){%
        $\begin{aligned}%
              z_f^T(k) = [ - & x_f(k-n_a),\,\cdots, -x_f(k-1),                       \\%
                             & \nu_f(k-n_b-n_k-1), \,\cdots, \,\nu_f(k-n_k) ]%
          \end{aligned}$%
    };

    % node système auxiliaire
    \node [block, below=\y of z.south, fill=teal!30, anchor=north] (Z) {%
        $Z_f = %
        \begin{bmatrix} %
            z_f^T(1) \\ %
            \vdots \\ %
            z_f^T(N) %
        \end{bmatrix}$%
        };

    % node theta_init
    \node [input] (theta_init) at ([xshift=-\x]r1-|Saux.west) {$\theta_\text{init}$};

    % node condition d'arrêt
    \node [block, fill=lightgray!50, anchor=north west] (stop) at ([xshift=-2*\x,yshift=-\y/2] y-|Fsvfr.east) {\scriptsize \text{while} $\displaystyle \sum_{j=1}^{\dim \hat{\theta}_{sriv}} \left|\frac{\hat{\theta}_{sriv,j}^{\text{iter}}-\hat{\theta}_{sriv,j}^{\text{iter-1}}}{\hat{\theta}_{sriv,j}^{\text{iter}}}\right| > \epsilon$};

    % ------------------------------
	% --- définition des flèches ---
	% ------------------------------
    % de r1 au sommateur
    \draw [->] (r1) |- (sum2);

    % de r2 au sommateur
    \draw [->] (r2) -- (CpCiPf);
    \draw [->] (CpCiPf) -- (sum2);

    % de y au sommateur
    \draw [->] (y) |- (CpCi);
    \draw [->] (CpCi) -- (comp2);

    % du sommateur au comparateur 
    \draw [->] (sum2) -- (comp2) node [midway, right] (r) {$r(k)$};

    % du comparateur au filtre 
    \draw [->] (comp2) -- (Fsvfl) node [midway, right] {$u(k)$};

    % de la sortie y au filtre
    \draw [->] (CpCi-|y) -- (Fsvfr);

    % des filtres à phi
    \draw [->] (Fsvfl) -- (phi.north-|Fsvfl) node [midway, right] {$u_f(k)$};
    \draw [->] (Fsvfr) -- (phi.north-|Fsvfr) node [midway, left, name=yf] {$y_f(k)$};

    % de phi à Phi
    \draw [->] (phi) -- (Phi);

    % de yf à Y
    \draw [->] (yf) -| (Y);

    % de Phi à la solution
    \draw [->] (Phi) -- (sol);

    % de Y à la solution
    \draw [->] (Y) |- (sol);

    % de sol à theta
    \draw [->] (sol) -- (theta) node [midway] (theta_s) {};

    % de r à Saux
    \draw [->] (r) -| ([yshift=\y]$(Saux.north east)!0.5!(CpCiPf.north west)$) -| (Saux);

    % de Saux à Gaux
    \draw [->] (Saux) -- (Gaux) node [midway, left] (nu) {$\nu(k)$};

    % de Gaux au filtre
    \draw [->] (Gaux) -- (Fsvflaux) node [midway, left] {$x(k)$};

    % de nu au filtre
    \draw [->] (nu) -| (Fsvfraux);

    % des filtres à z
    \draw [->] (Fsvflaux) -- (z.north-|Fsvflaux) node [midway, left] {$x_f(k)$};
    \draw [->] (Fsvfraux) -- (z.north-|Fsvfraux) node [midway, right] {$\nu_f(k)$};

    % de z à Z
    \draw [->] (z) -- (Z);

    % de z à sol
    \draw [->] (Z) |- (sol);

    % de theta_init à Saux et Gaux
    \draw [->, dotted] (theta_init) |- (Saux);
    \draw [->, dotted] (theta_init) |- (Gaux);

    % node iter+1
    \node [block, fill=lightgray!50] (iter) at ([xshift=-\dy]z.west|-theta_s) {$\text{iter} = \text{iter} + 1$};

    % de theta_s à iter
    \draw [->, densely dashed] (theta_s) -- (iter) node [very near end, above] {$\hat{\theta}_{sriv}^{\text{iter}}$};

    % de iter à Gaux/Saux et A^-1
    \begin{pgfonlayer}{background}
        \draw [->, densely dashed] (iter) |- ([xshift=\dy,yshift=-\y/3]Saux.south west) node [very near start, right] {$\hat{\theta}_{sriv}^{\text{iter}-1}$} -- ([xshift=-\dy,yshift=\y/3]Saux.north east);
        \draw [->, densely dashed] ([yshift=-\y/3]iter|-Gaux.south west) -- ([xshift=\dy,yshift=-\y/3]Gaux.south west) -- ([xshift=-\dy,yshift=\y/3]Gaux.north east);
        \draw [->, densely dashed] ([yshift=-\y/3]iter|-Fsvflaux.south west) -- ([xshift=\dy,yshift=-\y/3]Fsvflaux.south west) node [at end] (fleche_Fsvflaux) {} -- ([xshift=-\dy,yshift=\y/3]Fsvflaux.north east);
        \draw [->, densely dashed] (fleche_Fsvflaux) |- ([xshift=\dy,yshift=-\y/3]Fsvfraux.south west) node [at end] (fleche_Fsvfraux) {} -- ([xshift=-\dy,yshift=\y/3]Fsvfraux.north east);
        \draw [->, densely dashed] (fleche_Fsvfraux) |- ([xshift=\dy,yshift=-\y/3]Fsvfl.south west) node [at end] (fleche_Fsvfl) {} -- ([xshift=-\dy,yshift=\y/3]Fsvfl.north east);
        \draw [->, densely dashed] (fleche_Fsvfl) |- ([xshift=\dy,yshift=-\y/3]Fsvfr.south west) -- ([xshift=-\dy,yshift=\y/3]Fsvfr.north east);
    \end{pgfonlayer}

    % rectangle while
    \begin{pgfonlayer}{background}
        \draw [loosely dashdotted] ([xshift=0.25,yshift=0.25]iter.south west) rectangle ([xshift=-0.25,yshift=-0.25]stop.north east);
    \end{pgfonlayer}
    
    % % partie génération du signal
    % \coordinate (p1) at ([xshift=\dx, yshift=-\dy]CpCi.south);
    % \coordinate (p2) at ([yshift=-\dy]CpCi.south-|stop.east);
    % \draw [dashed] (p1) -- (p2);
    % \node [align=center, anchor=west] (sig_gen) at ($(r1.north-|p1)!0.5!(p1)$) {Génération\\du\\signal};

    % % partie filtrage
    % \coordinate (p3) at ([yshift=\dy]p1|-phi.north);
    % \coordinate (p4) at ([yshift=\dy]p2|-phi.north);
    % \draw [dashed] (p3) -- (p4);
    % \coordinate (ptextfilt) at ($(p3)!0.5!(p1)$);
    % \node (filt) at (ptextfilt-|sig_gen) {Filtrage};

    % % partie régresseur
    % \coordinate (p5) at ([yshift=-\dy]p1|-Phi.south);
    % \coordinate (p6) at ([yshift=-\dy]p2|-Phi.south);
    % \draw [dashed] (p5) -- (p6);
    % \coordinate (ptextreg) at ($(p5)!0.5!(p3)$);
    % \node [align=center] (reg) at (ptextreg-|sig_gen) {Conception\\du\\regresseur};

    % % partie résolution
    % \coordinate (ptextresol) at ($(p5)!0.5!(theta)$);
    % \node [align=center] (resol) at (ptextresol-|sig_gen) {Résolution};
\end{tikzpicture}